% Pursuing Stacks (cotes 134-2 to 134-8): archive cote 134-7, site page 80 (PDF page 81), typescript page 522.
% typescript. Draft compared: ps6.tex:3979-4037.
% AI-transcribed by Claude Opus 5.5 (claude-opus-5-5, adaptive thinking, effort high), Message Batches API, batch ps4 (msgbatch_012Aoc1NuskATPsyBeKSYK7J), 2026-11-12; not yet checked by eye.
\scanpage{134-7/80}
\mspage{522}
venient to denote the category of the latter by
\[ (3) \qquad Hotab_0(k) \]
rather than $D.(Ab)(0\text{-conn})$,  the subscript 0 denoting the 0-connectedness
condition.  The reason for this change in notation is that,  according
to the choices made in the basic definitions,  the ``k-linear'' algebraic
interpretation of this category may vary still.  Strictly speaking,  in
the approach we have been following in terms of semisimplicial unipotent
bundles,  the category (3) cannot really be described as the subcategory
of 0-connected objects of $D.(Ab_k)$,  we have seen rather that in order to
define a functor $\overset{\hand{\sim}}{LH}.$ ,  we have to work with \xout{semisimplicial} chain comple-%
xes in $Pro(Ab_k)$  rather than \xout{\unsure{in}} in $(Ab_k)$; \xout{(}cf. section 11\over{\illegible{1}}{5}, where this
point is made rather forcefully. (Recall that taking projective limits
instead would'nt help,  because then the linearization functor $\overset{\hand{\sim}}{LH}.$
would no longer commute to ground ring extension ! ).  The trouble then
is that in order to define
\[ \overset{\hand{\sim}}{K} : \ Hotab_0(k) \longrightarrow Hot_0(k) \quad , \]
we \deleted{\unsure{are}}\inserted{are} obliged to enlarge accordingly the category of models used for
describing  $Hot_0(k)$,  namely take semisimplicial \uline{pro}-unipotent bundles ,
rather than just semisimplicial unipotent bundles.  This sudden invasion
of the picture by proobjects may appear forbidding  -  but maybe it will
appear less so,  or at any rate kind of natural and unescapable,  if we re
\hand{-}call that when it will come to working with homotopy types of general
topoi,  these are anyhow ``prohomotopy types'' rather,  and they can be des-%
cribed only \xout{in terms} by working systematically in terms of proob\over{\unsure{g}}{\unsure{j}}ects of
various categories.  We'll see \deleted{also}\inserted{however} that there may be a way out of the
``pro''-mess,  by using a slightly different,  more or less dual approach
towards the idea of ``unipotent bundles''.  Of course,  one may also think
of using the ``pointed linearization'' $L_{pt}$ for $LH.$ ,  instead of $L$ ,
\xout{\unsure{in defining LH LH (X}} but as was pointed out in section 117,  this
will lead to ``formal'' homology and cohomology invariants rather than to
``schematic'' ones,  and the relation between these and corresponding inva-%
riants for ordinary homotopy types will be looser still;  at any rate,
it is suited for \xout{defining} describing Postnikov dévissage in the set-up
of ``formal schematic'' homotopy types,\inserted{only} not for schematic homotopy types.
\deleted{But}\inserted{Still,} the theory of the former may have an interest in it's own right,  even
though it's relation to ordinary homotopy types is'nt so clear at present.
Thus we may \inserted{state} that there are \inserted{around} at present \deleted{around} three or four different can-%
didates for possibly fitting the daydreaming I want to write down.  \deleted{\unsure{With}}
% NOTE: the right edge of the page is cut at "if we re" and at the struck word after "write down."
\margin{80}
